% concatenated from order-convexity-corrigendum.tex
\documentclass[11pt]{amsart}

% started: Feb 25, 2019
% last (update: not so last) edit: Sep 6, 2019
% more edits: Jan 31 - Feb 2, 2020 (following Prof. Schep's comments).
% final (hopefully) edits: Feb 17, 2020 (reply to referee)

\usepackage{comment}

\newtheorem{theorem}{Theorem}[section]
\newtheorem{acknowledgement}{Acknowledgments}
\newtheorem{algo}[theorem]{Algorithm}
\newtheorem{axiom}[theorem]{Axiom}
\newtheorem{case}[theorem]{Case}
\newtheorem{claim}[theorem]{Claim}
\newtheorem{conclusion}[theorem]{Conclusion}
\newtheorem{condition}[theorem]{Condition}
\newtheorem{conjecture}{Conjecture}
\newtheorem{corollary}[theorem]{Corollary}
\newtheorem{criterion}[theorem]{Criterion}
\newtheorem{definition}[theorem]{Definition}
\newtheorem{lemma}[theorem]{Lemma}
\newtheorem{notation}[theorem]{Notation}
\newtheorem{problem}[theorem]{Problem}
\newtheorem{proposition}[theorem]{Proposition}
% \newtheorem{remark}[theorem]{Remark}
\newtheorem{solution}[theorem]{Solution}
\newtheorem{summary}[theorem]{Summary}
% \newtheorem{example}[theorem]{Example}

\newtheorem*{theorem*}{Theorem}
\newtheorem*{corollary*}{Corollary}

\theoremstyle{definition}
% \newtheorem{definition}[theorem]{Definition}
\newtheorem{remark}[theorem]{Remark}
\newtheorem{question}[theorem]{Question}
\newtheorem{example}[theorem]{Example}
% \newtheorem{numbered}[theorem]{}

%%%%%%%%%%%%%%%%%%%%%%%%%%%%%%%%%%%%%%%%%%%%%%%%%%%%%%%%%%%%%%%%%%%%%%%%%%%%%%%%%%%%%%%%%%%%%%%%%%%
% \usepackage{showkeys}
% \usepackage{amsmath,amssymb,amsthm}
% \usepackage[backref, colorlinks, bookmarks=true]{hyperref}
\usepackage[colorlinks, bookmarks=true]{hyperref}
\usepackage{color,graphicx,shortvrb}
% \usepackage[active]{srcltx} %SRC Specials for DVI search

\usepackage{enumerate}
\usepackage{amsmath}
\usepackage{amssymb}

%%%%%%%%%%%%%%%%%%%%%%%%%%%%%%%%%%%%%%%%%%%%%%%%%%%%%%%%%%%%%%%%%%%%%%%%%%%%%

\def\query#1{\setlength\marginparwidth{80pt}%
% \marginpar{\raggedright\fontsize{7.81}{9}\selectfont\itshape\hrule\smallskip
\marginpar{\raggedright\fontsize{10}{10}\selectfont\itshape
\hrule\smallskip
{\textcolor{red}{#1}}\par\smallskip\hrule}}
\def\removequeries{\def\query##1{}} 

% \def\E{{\widehat{E}}}
% \def\W{{\widehat{W}}}
% \def\F{{\widehat{F}}}
% \def\A{{\widehat{A}}}
% \def\T{{\widehat{T}}}
% \def\J#1#2#3{ \left\{ #1,#2,#3 \right\} }
% \def\RR{{\mathbb{R}}}
% \def\NN{{\mathbb{N}}}
% \def\11{\textbf{$1$}}
% \def\CC{{\mathbb{C}}}
% \def\HH{{\mathbb{H}}}
% \def\KK{{\mathbb{K}}}
% \def\C#1{ \mathcal{#1} }

%\usepackage[notcite,notref]{showkeys}

\newcommand{\ppe}{\hat{\otimes}_{\epsilon}}
\newcommand{\ppi}{\hat{\otimes}_{\pi}}
\newcommand{\lra}{\longrightarrow}
\newcommand{\w}{\omega}
\newcommand{\uno}{1\!\!1}
\newcommand{\supp}{\mathrm{supp}}
\newcommand{\is}{\mathcal{S}}
\newcommand{\boa}{\boldsymbol{\alpha}}
\newcommand{\bob}{\boldsymbol{\beta}}
\newcommand{\tr}{\mathrm{tr}}
\newcommand{\ce}{\mathcal{E}}
\newcommand{\se}{\is_{\mathcal{E}}}
\newcommand{\mx}{\mathrm{max}}
\newcommand{\mn}{\mathrm{min}}
\newcommand{\pip}{\overline{\Pi}_2}
% \newcommand{\HH}{\mathbb{H}}

\newcommand{\torus}{\mathbb{T}}
\newcommand{\tri}[1]{| \! | \! | {#1} | \! | \! |}
\newcommand{\one}{\mathbf{1}}
\newcommand{\pp}{\mathfrak{P}}
\newcommand{\ran}{\mathrm{ran} \,}
\newcommand{\rank}{\mathrm{rank}}
\newcommand{\spn}{\mathrm{span}}
\newcommand{\vr}{\varepsilon}
\newcommand{\ave}{{\mathrm{Ave}}}
\newcommand{\diag}{{\mathrm{diag}}}
\newcommand{\DP}{\mathrm{DP}}
\newcommand{\BP}{\mathrm{BP}}
\newcommand{\ball}{\mathbf{B}}
\newcommand{\sphere}{\mathbf{S}}
\newcommand{\ep}{\mathrm{EP}}
\newcommand{\oep}{\mathrm{OEP}}
\newcommand{\expe}{\mathbb{E}}
\newcommand{\prob}{\mathbb{P}}
\newcommand{\tensa}{\otimes_\alpha}

\newcommand{\N}{\mathbb{N}}
\newcommand{\Z}{\mathbb{Z}}
\newcommand{\C}{\mathbb{C}}
\newcommand{\R}{\mathbb{R}}
\newcommand{\Q}{\mathbb{Q}}

\newcommand{\A}{\mathcal{A}}
\newcommand{\ato}{\mathfrak{A}}

\def \eqalign#1{\null\,\vcenter{\openup\jot %%%   \m@th
   \ialign{\strut\hfil$\displaystyle{##}$&$
      \displaystyle{{}##}$\hfil \crcr#1\crcr}}\,}


\newcommand{\pprec}{\prec \!\! \prec}
\newcommand{\dist}{\mathrm{dist}}
% \newcommand{\injtens}{\check{\otimes}}
% \newcommand{\projtens}{\hat{\otimes}}
% \newcommand{\anytens}{\tilde{\otimes}}
% \newcommand{\kwap}{{\mathrm{Kw}}}
% \newcommand{\col}{{\mathbf{C}}}
% \newcommand{\seq}{{\mathbf{s}}}
\newcommand{\csch}{{\mathrm{CSCH}}}
\newcommand{\sch}{{\mathrm{SCH}}}
\newcommand{\ch}{{\mathrm{CH}}}
\newcommand{\so}{{\mathrm{S}}}

\newcommand{\triple}[1]{{\left\vert\kern-0.25ex\left\vert\kern-0.25ex\left\vert #1 
    \right\vert\kern-0.25ex\right\vert\kern-0.25ex\right\vert}}

\begin{document}

\numberwithin{equation}{section}

% \title[Analogues of Krein-Milman]{Order extreme points, solid convex hull, and lattice analogues of Krein-Milman Theorem}

\title[Order extreme points]{Corrigendum: order extreme points \\ and solid convex hulls}

% \author{T.~Oikhberg and M.A.~Tursi -- REDO, ADD Anastasiia Ianina}

% \address{ Dept.~of Mathematics, University of Illinois, Urbana IL 61801, USA}
% \email{oikhberg@illinois.edu, gramcko2@illinois.edu}

\author{A.~Ianina and T.~Oikhberg}

\address{A.I. and T.O.: Dept.~of Mathematics, University of Illinois, Urbana IL 61801, USA}
\email{aianina2@illinois.edu, oikhberg@illinois.edu}

 \author{M.A.~Tursi}


\address{M.A.T.: Independent researcher}
\email{maryangelica.tursi@gmail.com}

\date{\today}

\subjclass[2010]{46B22, 46B42}

\keywords{Banach lattice, extreme point, convex hull, Radon-Nikod{\'y}m Property}

% \dedicatory{To the memory of Victor Lomonosov}

% \thanks{We wish to thank the organizers of Positivity IX, where part of this work was carried out.}

% \date{}
\maketitle

\parindent=0pt
\parskip=3pt

\begin{abstract}
We correct some errors found in [T.~Oikhberg and M.A.~Tursi,
{\it Order extreme points and solid convex hulls},
The Mathematical Legacy of Victor Lomonosov (ed.~R.~Aron et.al.), de Gryuter, 2020, 297--315.]
\end{abstract}


\maketitle
\thispagestyle{empty}

% \section{Introduction}\label{s:intro}

The goal of this note is two-fold: (i) to correct some annoying typos in \cite{OiTu}, and (ii) to correct some of the proofs. The enumeration of sections and statements (theorems, definitions, examples etc.) in the Arxiv preprint is different from that in the published version; the numbers from the published version will be given in square brackets.


{\bf (1)} In several statements in Sections 3 [19.3] and 4 [19.4], the convexity assumptions are omitted. Specifically:

$\bullet$ Proposition 3.1 [19.7] -- Separation -- should read: \\
Suppose $\tau$ is a sufficiently rich topology on a Banach lattice $X$, and $A \subset X_+$ is a $\tau$-closed positive-solid bounded {\bf convex} subset of $X_+$. Suppose, furthermore, $x \in X_+$ does not belong to $A$. Then there exists $f \in X^\tau_+$ so that $f(x) > \sup_{a \in A} f(a)$.

$\bullet$ Theorem 4.1 [19.10] -- ``Solid'' Krein-Milman -- should read: \\
Any $\tau$-compact positive-solid {\bf convex} subset $A$ of $X_+$ coincides with the $\tau$-closed positive-solid convex hull of its order extreme points.

$\bullet$ Corollary 4.2 [19.11] should read: \\
Any $\tau$-compact solid {\bf convex} subset of $X$ coincides with the $\tau$-closed solid convex hull of its order extreme points.

Without convexity, these results may fail. For instance, for Corollary 4.2, we can take $X = (\R^2, \| \cdot \|_\infty)$, and $A = \{(x,0) : |x| \leq 1\} \cup \{(0,y) : |y| \leq 1\}$. The set $A$ is closed when $\tau$ is the norm topology and solid, but the convex hull of $A$ is much larger than $A$ itself.

{\bf (2)} There is a hole in the proof of Theorem 7.1 [19.30]: on page 17 line -2 [page 313 line -6], it is claimed that, for any $x \in C_+ \backslash \{0\}$, we have $x \wedge u \neq 0$. The current argument doesn't permit to reach this conclusion. However, this can be fixed as follows:

\begin{proof}[Patch on the proof of Theorem 19.30.]
Denote by $B \subset X$ the band generated by $u$ in $X$. By \cite[Proposition 1.2.3]{M-N}, $B$ is closed (and solid); by \cite[Proposition 1.2.3]{M-N}, $B_+$ consists of all $z \in X_+$ so that $z = \vee_n (z \wedge nu)$.
As $u$ is a quasi-interior point of $Y$, we conclude that $C' \subset Y \subset B$, hence $C \subset B$. So, $x = \vee_n (x \wedge nu)$, which implies $x \wedge nu \neq 0$ for large enough $n$. As $x \wedge nu \leq n (x \wedge u)$, we conclude that $x \wedge u \neq 0$.
%
% \cite[Section 1.2]{M-N} re bands and band projections.
% By \cite[Page 18]{M-N}, $B$ is a projection band, although we do not seem to need that. 
\end{proof}



{\bf (3)} The proof of the implication $(1) \Rightarrow (2)$ in Proposition 4.3 [19.12] is incorrect: the set $D$ constructed on page 10 [306], first paragraph, is not solid, hence it is not guaranteed to have any order extreme points.

Proposition 4.3 [19.12] can be established in a rather roundabout way. More specifically, Theorem 7.1 [19.30] implies:

\begin{corollary*}
For a Banach lattice $X$, the following statements are equivalent:
\begin{enumerate}
	\item $X$ has the Radon-Nikodym Property $($RNP$)$.
	\item $X$ has the Krein-Milman Property $($KMP$)$.
	\item $X$ has the Solid Krein-Milman Property $($SKMP$)$ -- that is, any closed bounded convex solid subset of $X$ is the closure of the solid convex hull of its order extreme points.
	\item Any closed bounded convex solid subset of $X$ has an order extreme points.
\end{enumerate}
\end{corollary*}

Above, item (4) is a \emph{verbatim} repetition of Proposition 4.3[19.12](1), while item (3) is nothing but a restatement of Proposition 4.3[19.12](2).

\begin{proof}
The implication $(3) \Rightarrow (4)$ is immediate, and $(2) \Rightarrow (3)$ follows by combining Theorem 2.2 [19.2] with Corollary 2.4 [19.4]. The equivalence $(1) \Leftrightarrow (2)$ is known, see the references cited in \cite{OiTu}.

The proof of  Theorem 7.1 [19.30] actually produces, for any Banach lattice $X$ failing the RNP, a closed bounded convex solid $E \subset X$ without any order extreme points. This gives $\lnot (1) \Rightarrow \lnot (4)$, or equivalently, $(4) \Rightarrow (1)$.
\end{proof}


\begin{thebibliography}{22}

\bibitem{M-N} P. Meyer-Nieberg.
Banach lattices. Springer, Berlin, 1991.

\bibitem{OiTu} T.~Oikhberg and M.A.~Tursi,
Order extreme points and solid convex hulls.
The Mathematical Legacy of Victor Lomonosov (ed.~R.~Aron et.al.), de Gryuter, 2020, 297--315.
\url{https://arxiv.org/abs/1907.00660}


\end{thebibliography}

\end{document}

% \bibitem{Ab} Y.~Abramovich.
% Isometries of normed lattices (in Russian).
% Optimizatsiya 43(60) (1988), 74--80. 

%\bibitem{AK} F. Albiac and N. Kalton. 
%Topics in Banach space theory, second edition.
%Springer, Cham, 2016.

\bibitem{AB} C. Aliprantis and O. Burkinshaw.
Positive operators.
Springer, Berlin, 2006.

\bibitem{Ama} M. Amarante.
The Sandwich theorem via Pataraia’s fixed point theorem.
Positivity 23 (2019), 97--100.

\bibitem{ArL} R. Aron and V. Lomonosov.
After the Bishop-Phelps theorem.
Acta Comment. Univ. Tartu. Math. 18 (2014), 39--49. 

% \bibitem{Be} S. Bellenot. Banach spaces with trivial isometries.
% Israel J. Math. 56 (1986), 89--96. 

%\bibitem{BeP} C. Bessaga and A. Pelczynski.
%On extreme points in separable conjugate spaces.
%Israel J. Math. 4 (1966), 262--264. 

%\bibitem{BP} E. Bishop and R. Phelps.
%The support functionals of a convex set.
%1963 Proc. Sympos. Pure Math., Vol. VII pp. 27--35 Amer. Math. Soc., Providence, R.I.

\bibitem{Boll} B. Bollob\'as.
An extension to the theorem of Bishop and Phelps.
Bull. London Math. Soc. 2 (1970), 181--182. 

\bibitem{BourTal} J.  Bourgain and  M.  Talagrand.
Dans un  espace de  Banach r\'eticul\'e solide, la  propri\'et\'e de  Radon-Nikodym et  celle de  Krein-Milman sont equivalentes. Proc. Amer. Math. Soc. 81  (1981), 93--96.

\bibitem{Bour} R. Bourgin.
Geometric aspects of convex sets with the Radon-Nikod{\'y}m property.
% Lecture Notes in Mathematics, 993.
Springer, Berlin, 1983.

\bibitem{Cas} V. Caselles.
A short proof of the equivalence of KMP and RNP in Banach lattices and preduals of von Neumann algebras. Proc. Amer. Math. Soc. 102 (1988), 973-–974. 

% \bibitem{CaGo} J. Castillo and M. Gonzalez.
% Three-space problems in Banach space theory.
% Lecture Notes in Mathematics, 1667. 
% Springer, Berlin, 1997.

\bibitem{ChW} Z.L. Chen and A. Wickstead.
Relative weak compactness of solid hulls in Banach lattices.
Indag. Math. (N.S.) 9 (1998), 187--196. 

\bibitem{CKMetc}
M. Chica, V. Kadets, M. Mart\'in, S. Moreno-Pulido, and F. Rambla-Barreno.
Bishop-Phelps-Bollob\'as moduli of a Banach space.
J. Math. Anal. Appl. 412 (2014), 697--719. 

% \bibitem{Cow} E.R. Cowie.
% A note on uniquely maximal Banach spaces.
% Proc. Edinburgh Math. Soc. (2) 26 (1983), 85--87. 

% \bibitem{DGL} W. Davis, N. Ghoussoub, and J. Lindenstrauss.
% A lattice renorming theorem and applications to vector-valued processes.
% Trans. Amer. Math. Soc. 263 (1981), 531--540. 

% \bibitem{DLOT} H.~G.~Dales, N.~Laustsen, T.~Oikhberg, and V.~Troitsky.
% Multi-norms and Banach lattices.
% Dissertationes Math. (Rozprawy Mat.) 524 (2017). % , 115 pp..

\bibitem{DU} J. Diestel and J. Uhl.
Vector measures.  % With a foreword by B. J. Pettis. Mathematical Surveys, No. 15.
American Mathematical Society, Providence RI, 1977. % xiii+322 pp. 

\bibitem{EWe} E. Effros and C. Webster.
Operator analogues of locally convex spaces, Operator algebras and applications.
% (Samos, 1996), 
NATO Adv. Sci. Inst. Ser. C Math. Phys. Sci., vol. 495, Kluwer Acad. Publ., Dordrecht, 1997, 163--207. 

\bibitem{EWi} E. Effros and S. Winkler.
Matrix convexity: operator analogues of the bipolar and Hahn-Banach theorems.
J. Funct. Anal. 144 (1997), 117--152. 

\bibitem{FHHMZ} M. Fabian, P. Hajek, P. Habala, V. Montecinos, V. Zizler.
Banach space theory. The basis for linear and nonlinear analysis.
% CMS Books in Mathematics/Ouvrages de Mathématiques de la SMC.
Springer, New York, 2011.

\bibitem{FaM} D. Farenick and P. Morenz.
$C^*$-extreme points in the generalised state spaces of a $C^*$-algebra. 
Trans. Amer. Math. Soc. 349 (1997), 1725--1748.

% \bibitem{FeG_TAMS10} V. Ferenczi and E.M. Galego.
% Countable groups of isometries on Banach spaces.
% Trans. Amer. Math. Soc. 362 (2010), 4385--4431. 

% \bibitem{FeR_EM11} V. Ferenczi and C. Rosendal.
% Displaying Polish groups on separable Banach spaces.
% Extracta Math. 26 (2011), 195--233. 

% \bibitem{FeR_DU13} V. Ferenczi and C. Rosendal.
% On isometry groups and maximal symmetry.
% Duke Math. J. 162 (2013), 1771--1831. 

%\bibitem{Fo} V. Fonf.
%A property of Lindenstrauss-Phelps spaces.
%Functional Anal. Appl. 13 (1979), 66--67.

\bibitem{FW} B. Fuchssteiner and J.D. Maitland Wright.
Representing isotone operators on cones. Q. J. Math. 28 (1977), 155–-162.

% \bibitem{GTWZ} G. Gorefroy, S. Troyanski, J. Whitfield, and V. Zizler.
% Three-space problem for locally uniformly rotund renormings of Banach spaces.
% Proc. Amer. Math. Soc. 94 (1985), 647--652. 

% \bibitem{HKM} H. Hudzik, A. Kami{\'n}ska, and M. Masty{\l}o.
% Monotonicity and rotundity properties in Banach lattices.
% Rocky Mountain J. Math. 30 (2000), 933--950. 

%\bibitem{Kad} V. Kadets.
%A course in functional analysis and measure theory.
%% Translated from the 2006 Russian edition [MR2268285] by Andrei Iacob. Universitext. 
%Springer, Cham, 2018.

% \bibitem{Ke} A.~Kechris.
% Classical descriptive set theory.
% Springer, Berlin, 1995.

%\bibitem{LeT} D. Leung and W.K. Tang.
%Persistence of Banach lattices under nonlinear order isomorphisms, Positivity, 20 (2016), 709--717.

\bibitem{Linl1} J. Lindenstrauss.
On extreme points in $\ell_1$.
Israel J. Math. 4 (1966), 59--61. 

\bibitem{LP} J. Lindenstrauss and R. Phelps.
Extreme point properties of convex bodies in reflexive Banach spaces.
Israel J. Math. 6 (1968), 39--48. 

%\bibitem{LT1} J. Lindenstrauss and L. Tzafriri.
%\newblock {\em Classical Banach spaces I},
%\newblock Springer, Berlin, 1977.

% \bibitem{LT2} J. Lindenstrauss and L. Tzafriri.
% \newblock {\em Classical Banach spaces II},
% \newblock Springer, Berlin, 1979.

%\bibitem{Lo} H. Lotz.
%Extensions and liftings of positive linear mappings on Banach lattices.
%Trans. Amer. Math. Soc. 211 (1975), 85--100. 

\bibitem{Mag} B. Magajna.
$C^*$-convex sets and completely positive maps.
Integral Equations Operator Theory 85 (2016), 37--62. 

% \bibitem{Ma} B. Maurey.
% Symmetric distortion in $\ell_2$. Geometric aspects of functional analysis (Israel, 1992--1994), 143--147,
% Oper. Theory Adv. Appl., 77, Birkh{\"a}user, Basel, 1995. 

\bibitem{M-N} P. Meyer-Nieberg.
Banach lattices. Springer, Berlin, 1991.

% \bibitem{OdSchAM} E. Odell and T. Schlumprecht.
% The distortion problem. Acta Math. 173 (1994), 259--281. 

% \bibitem{OdSchHB} E. Odell and T. Schlumprecht.
% Distortion and asymptotic structure. Handbook of the geometry of Banach spaces, Vol. 2, 1333--1360, North-Holland, Amsterdam, 2003. 

% \bibitem{Oi} T. Oikhberg.
% A note on latticeability and algebrability. J. Math. Anal. Appl. 434 (2016), 523--537. 

\bibitem{PR} M. Popov and B. Randrianantoanina.
Narrow operators on function spaces and vector lattices.
% De Gruyter Studies in Mathematics, 45. 
Walter de Gruyter, Berlin, 2013.

% \bibitem{Re} C. Read.
% When $E$ and $E[E]$ are isomorphic. Geometry of Banach spaces (Strobl, 1989), 245--252,
% % London Math. Soc. Lecture Note Ser., 158, 
% Cambridge Univ. Press, Cambridge, 1990. 

\bibitem{Rud} W. Rudin.
Functional analysis.
McGraw-Hill, Singapore, 1991.

\bibitem{Sch} H.H. Schaefer.
Banach lattices and positive operators.
% Die Grundlehren der mathematischen Wissenschaften, Band 215. 
Springer, New York--Heidelberg, 1974.

\bibitem{WeWi} C. Webster and S. Winkler.
The Krein-Milman theorem in operator convexity.
Trans. Amer. Math. Soc. 351 (1999), 307--322. 

% \bibitem{Woo} G. Wood.
% Maximal algebra norms. 
% Trends in Banach spaces and operator theory (Memphis, TN, 2001), 335--345,
% Contemp. Math., 321, Amer. Math. Soc., Providence, RI, 2003. 

\end{thebibliography}

\end{document}

=========================================================================================

ON SHADOWS

The results above use the assumption that $A_s$ is closed. The following two lemmas show that, under certain conditions, $A_s$ is indeed closed.

\begin{lemma}\label{l:shadow_compact}
If $A \subset X_+$ is compact, then $A_s$ is closed.
\end{lemma}

\begin{proof}
Suppose the net $(x_\alpha)_{\alpha \in {\mathcal{I}}} \subset A_s$ converges to $x$, and show that $x \in A_s$ as well.
By the definition of $A_s$, for any $\alpha$ there exists $a_\alpha \in A$ so that $x_\alpha \leq a_\alpha$.
By compactness, the net ${\mathcal{I}}$ has a cofinite subnet ${\mathcal{J}}$ so that $(a_\beta)_{\beta \in {\mathcal{J}}}$ converges to some $a \in A$. Then $a \geq x$, hence $x \in A_s$. 
\end{proof}


ON WEAK* COMPACTNESS:

The following also immediately follows from Lemma \ref{l:shadow_compact} and the fact that the positive cone is weak$^*$-closed.

\query{This corollary is never used I think; delete?}

\begin{corollary}\label{c:switch-close-solid2}
	Suppose $X$ is a dual Banach lattice, and $C\subseteq X_+$ is bounded.  Then $S(\overline{C}^*)$ is weakly closed.
\end{corollary}


========================================================================================

PROOF OF SKMP <=> RNP


	If $X$ has the KMP, then if $C$ is a closed, bounded solid subset of $X_+$, then $S(C)$ is closed, and $C=S(C)_+$.  From the proof in theorem \ref{t:connection}, $OE(C) = OE(S(C)) = EP(S(C))_+$.  Since $X$ has the $KMP$, $S(C)$ is the closed convex hull of $EP(S(C))$. It follows that $C$ is the closed order convex hull of $OE(C)$.
	
	Suppose $X$ does not have the $KMP$. We now consider two cases.
	
	If $X$ is not $KB$, then there exists an sublattice $Y$ isomorphic to $c_0$ Let $C = \overline{S(B(Y)_+)}_+$. $C$ is closed,bounded, positive-solid in $X$, and convex.  We now show that $C$ has no order extreme points. Suppose $x\in OE(C)$.  Let $x_n \rightarrow x$, $x_n \leq y_n$, where $y_n\in B(Y)_+$.  Let $e_i$ be the atoms of $Y$.  We can without loss of generality we assume that each $x_n \leq y_n$, where $y_n$ is supported by finitely many atoms.  Let $\varepsilon$ be sufficiently small, and choose $n$ such that $\|x_n - x\| < \varepsilon$.  Choose $k$ such that $x_n \leq y_n \leq \sum_{i=1}^{k-1} e_i$.  Then $\| x\wedge e_k \| < \varepsilon$, Thus $x \vee e_k \gneq x$.  Furthermore, $x_n \vee e_k \leq y_n \vee e_k \in B(Y)$, and $x_n \vee e_k \rightarrow x \vee e_k$, so $ x \vee e_k \in C$.  Therefore, since $x \lneq x \vee e_k$, $x$ cannot be order extreme. 
	
	If $X$ is $KB$, but does not have the KMP, $X$ is an ideal of $X^{**}, $ so by the proof in Theorem 1 of \cite{Cas} there exists a closed, bounded, convex set $C\subseteq X_+$ such that $C$ has no extreme points, hence it has no order extreme points.  By proposition \ref{p:conv_closed_KB}, $S(C)_+$ is norm closed, bounded, positive-solid, and convex  but it has no extreme points since it is dominated by elements in $C$. Thus $X$ does not have the SKMP.


=========================================================================================

\bibitem{Ab} Y. Abramovich.
Isometries of normed lattices (in Russian).
Optimizatsiya 43(60) (1988), 74--80. 


\bibitem{BY} I. Ben Yaacov.
Fra\"iss\'e limits of metric structures.
J. Symb. Log. 80 (2015), 100--115. 

\bibitem{DLOT} H.~G.~Dales, N.~Laustsen, T.~Oikhberg, and V.~Troitsky.
Multi-norms and Banach lattices.
Dissertationes Math. (Rozprawy Mat.), to appear.

\bibitem{Do} K. Donner.
Extension of positive operators and Korovkin theorems.
% Lecture Notes in Mathematics, 904. 
Springer, Berlin-New York, 1982. % xii+182 pp. 

\bibitem{Lo} H.~Lotz.
Extensions and liftings of positive linear mappings on Banach lattices.
Trans. Amer. Math. Soc. 211 (1975), 85--100.

\bibitem{Lu} M. Lupini. ????

\bibitem{MaN} J. Marcolino Nhani.
La structure des sous-espaces de treillis.
Dissertationes Math. (Rozprawy Mat.) 397 (2001). % 50 pp..

\bibitem{McC} L.~B.~McClaran.
Some results on lattice structures on Banach spaces.
Ph.D. Dissertation, Texas A\&M University, 1994.

\bibitem{Ble} D. Blecher.
The standard dual of an operator space.
Pacific J. Math. 153 (1992) 15--30. 

\bibitem{CN} P. Casazza and N. Nielsen.
Embeddings of Banach spaces into Banach lattices and the Gordon-Lewis property.
Positivity 5 (2001) 297--321. 

\bibitem{DLOT} H.~G.~Dales, N.~Laustsen, T.~Oikhberg, and V.~Troitsky.
Multi-norms and Banach lattices.
Dissertationes Math. (Rozprawy Mat.), to appear.

\bibitem{DP} H.~G.~Dales and M.~E.~Polyakov.
Homological properties of modules over group algebras.
Proc. London Math. Soc. (3) 89 (2004) 390--426. 

% \bibitem{Daws} M. Daws.
% $p$-operator spaces and Fig\`a-Talamanca-Herz algebras.
% J. Operator Theory 63 (2010) 47--83. 

\bibitem{DFS} J. Diestel, J. Fourie, and J. Swart.
The metric theory of tensor products.
Grothendieck's r\'esum\'e revisited. American Math. Society, Providence, RI, 2008.

% \bibitem{DJT} J. Diestel, H. Jachrow, and A. Tonge.
% % % Diestel, Joe; Jarchow, Hans; Tonge, A.
% % {\em Absolutely summing operators}, % Cambridge Studies in Advanced Mathematics, 43.
% Absolutely summing operators,
% Cambridge University Press, Cambridge, UK, 1995.

\bibitem{DU} J. Diestel and J.~Uhl.
% {\em Vector measures}, % Cambridge Studies in Advanced Mathematics, 43.
Vector measures.
American Math. Society, Providence, RI, 1977.

\bibitem{ER} E. Effros and Z.-J. Ruan.
% {\em Operator spaces}, % Cambridge Studies in Advanced Mathematics, 43.
Operator spaces. Oxford University Press, Oxford, UK, 2000.

% \bibitem{LeM} C. Le Merdy.
% Factorization of $p$-completely bounded multilinear maps,
% Pacific J. Math. 172 (1996) 187--213.

% \bibitem{LT2} J. Lindenstrauss and L. Tzafriri.
% \newblock {\em Classical Banach spaces II},
% \newblock Springer, Berlin, 1977.

% \bibitem{Lam} A. Lambert.
% Operatorfolgenr\"aume, Eine Kategorie auf dem Weg von den Banach-Raumen zu den Operatorr\"aumen.
% Dissertation zur Erlangung des Grades Doktor der Naturwissenschaften,
% Universit\"at des Saarlandes, 2002.


% \bibitem{LNR} A. Lambert, M. Neufang, and V. Runde.
% Operator space structure and amenability for Figa-Talamanca-Herz algebras,
% J. Funct. Anal. 211 (2004) 245--269.

% \bibitem{LO} R.~Latala and K.~Oleszkiewicz.
% On the best constant in the Khinchin-Kahane inequality.
% Studia Math. 109 (1994), 101--104. 

% \bibitem{Lo} H.~Lotz.
% Extensions and liftings of positive linear mappings on Banach lattices.
% Trans. Amer. Math. Soc. 211 (1975), 85--100.

% \bibitem{Lu} M. Lupini.
% Convex $p$-multinormed spaces are mutiisometric to $p$-operator spaces,
% in preparation.

\bibitem{MaN} J. Marcolino Nhani.
La structure des sous-espaces de treillis.
Dissertationes Math. (Rozprawy Mat.) 397 (2001). % 50 pp..

\bibitem{McC} L.~B.~McClaran.
Some results on lattice structures on Banach spaces.
Ph.D. Dissertation, Texas A\&M University, 1994.

\bibitem{M-N} P. Meyer-Nieberg.
Banach lattices. Springer, Berlin, 1991.

\bibitem{MS}
V. Milman and G. Schechtman.
Asymptotic theory of finite-dimensional normed spaces. % LMN 1200
Springer, Berlin, 1986.

% \bibitem{NS} N. Nemesh and S. Shteyner.
% Metric and topological freedom for operator sequence spaces,
% in preparation.

% \bibitem{OiLR} T. Oikhberg.
% A note on local theory of $p$-multinormed spaces,
% in preparation.

\bibitem{OiRi} T. Oikhberg and E. Ricard.
Operator spaces with few completely bounded maps.
Math. Ann. 328 (2004), 229--259. 

\bibitem{PR} A. Pelczynski and H. Rosenthal.
Localization techniques in $L_p$ spaces.
% {\em Studia Math.}
Studia Math. 52 (1974/75), 263--289.

% \bibitem{PiINT} G.~Pisier.
% Complex interpolation and regular operators between Banach lattices.
% Arch. Math. (Basel) 62 (1994), 261--269.

% \bibitem{PiFACT} G.~Pisier.
% Factorization of linear operators and geometry of Banach spaces.
% % CBMS Regional Conference Series in Mathematics, 60.
% % Published for the Conference Board of the Mathematical Sciences, Washington, DC; by the
% American Mathematical Society, Providence, RI, 1986.

\bibitem{Rya} R.~Ryan.
Introduction to tensor products of Banach spaces. 
% Springer Monographs in Mathematics. 
Springer, London, 2002.

% \bibitem{Sz} S.~Szarek.
% On the best constants in the Khinchin inequality.
% Studia Math. 58 (1976), 197--208. 

% \bibitem{Tak} M. Takesaki.
% Theory of Operator Algebras I,
% Springer, 2002.

% \bibitem{T-J} N. Tomczak-Jaegermann,
% % \newblock {\em Banach-Mazur distances and finite-dimensional operator ideals},
% Banach-Mazur distances and finite-dimensional operator ideals,
% \newblock Longman Scientific and Technical, 1989.

=======================================================================================


We first describe a procedure for ``atomization'' of a separable $C(K)$.
We want to embeds $C(K)$ (lattice isometrically) into a separable atomic lattice $C(K^\rho)$.

Recall that, for a Hausdorff compact $K$, the space $C(K)$ is separable iff
$K$ metrizable. Henceforth we assume that $K$ is equipped with the metric $d$,
and has diameter $1$.

Inside of $K$, find a countable dense subset $\{k_1, k_2, \ldots \}$.
% We now equip $K$ with the distance $\rho$, defined as follows:
Define $\rho : K \times K \to [0,1]$ as follows:
\begin{itemize}
\item 
If $x, y \notin \{k_1, k_2, \ldots \}$, or if $x=y$, then $\rho(x,y) = d(x,y)$.
\item
If $x \notin \{k_1, k_2, \ldots \}$ and $y = k_i$ for some $i$, then
$\rho(x,y) = \rho(y,x) = \max\big\{d(x,y), 2^{-i}\big\}$.
\item
If $x = k_i$ and $y = k_j$, with $i > j$, then
$\rho(x,y) = \rho(y,x) = 2^{-j}$. % $\max\big\{d(x,y), 2^{-i}\big\}$.
\end{itemize}
One can check that $\rho$ is indeed a metric.
We denote by $K^\rho$ the set $K$ equipped with the topology described by $\rho$
(as opposed to the original topology generated by $d$).
One can check that open sets in the $\rho$-topology are of the form
$U \cup S$, where $U$ is open in the $d$-topology, and $S \subset \{k_1, k_2, \ldots \}$.
Further, the fact that any $f \in C(K^\rho)$ is $\rho$-uniformly continuous implies that
$f \in C(K^\rho)$ iff there exists $g \in C(K)$ so that $f|_{K \backslash \{k_1, k_2, \ldots \}} = g$.
Consequently, the atoms of $C(K^\rho)$ are characteristic functions of $\delta_{k_i}$.
Thus, $C(K^\rho)$ is atomic.


===============================================================================


THIS WORKS NOT ONLY FOR LATTICES BUT ALSO FOR GENERAL ORDERED SPACES

\begin{proposition}[Separation]\label{p:separation}
Suppose $A \subset X_+$ is a $\tau$-closed positive-solid bounded subset of $X_+$ so that $A_s$ is $\tau$-closed. Suppose, furthermore, $x \in X_+$ does not belong to $A$. Then there exists $f \in X^\tau_+$ so that $f(x) > \sup_{a \in A} f(a)$.
\end{proposition}

\begin{proof}
By Hahn-Banach Theorem, we can find $f \in X^\tau$ strongly separating $x$ from $A_s$: $f(x) > \sup_{a \in A} f(a)$. But $f \geq 0$. Indeed, otherwise, find $y \geq 0$ with $f(y) < 0$. Then $Ny \in A_s$ for any $N$, but $f(Ny) > f(x)$ for $N$ large enough.
\end{proof}

